\unNumberedSection{Laboratory Exercise}
Some general comments is that it was a good lab and I feel like I learnt something. The setup was easy to work with, and I liked the approach of letting us think a bit on what to measure and how. 

In the first part of the lab we aligned the laser in steps. We started by using an alignment laser and aligning the beam splitter and then detector B. Detectors T and A were already aligned. We then used the actual diode laser and aligned the BBO-crystal. The first experiment was to test if the light coming out of from the down-conversion is photon pairs. This was tested by measuring the coherence function for detectors T and B. It was found that $\qqcoherence_{TB} \approx 120 > 1$, so the light was very correlated, and we have a photon pair source. For the second experiment, single photon, we first needed to realize that the down-conversion creates thermal photons, so we cannot measure $\qqcoherence_{AB}$ directly, but by measuring $\qqcoherence_{TAB}$ we can make sure to only count the coincidences for detectors A and B when we also detect in detector T. Thus, by measuring in T we collapse, or project, the state unto the number state basis. The highest probability state will be $\ket{1}$, and we therefore suspect a coherence function close to zero which we measure. A detail is that we cannot measure the true coherence function for thermal light since the coherence time is much faster than the resolution of our detectors. Lastly we did an experiment by replacing the crystal with a fluorescent filter, which once again produced thermal light.

Realizing that the down-conversion was thermal light was a bit tricky, but discussing with the other students for a while led us to that conclusion. Thus, when tracing out the signal (idler) beam we get thermal photon statistics on the idler (signal) beam. I had also not realized how extremely short the coherence time is for thermal light. When we understood how the down-conversion worked, it was relatively easy to figure out how to perform the experiment, which then yielded the expected result. We also got an explanation of how fluorescence works and why it also gives thermal light. \study{I would really enjoy working more with the setup and maybe building more instead of just aligning. It would also be interesting to build a Bell experiment.}